\chapter{Curvas}\label{ch:b2:curves}

A geometria agora se torna diferencial. Uma curva é um ponto que se move pelo
espaço; o cálculo nos dá sua velocidade e sua aceleração, e a geometria
pergunta o que é intrínseco --- independente da rapidez com que percorremos a
trajetória. As respostas são o \emph{\omterm{def:b2:curves:length}{comprimento de arco}}, que mede a
própria trajetória, e a \emph{\omterm{thm:b2:curves:frenet2d}{curvatura}}, que mede quanto ela
se dobra. Em dimensão $3$ um segundo invariante, a \emph{\omterm{thm:b2:curves:frenet3d}{torção}},
mede quanto a curva se torce para fora de seu plano. O dispositivo de
contabilidade de tudo isso é o \emph{referencial de Frenet} móvel.

\section{Arcos parametrizados}

\begin{definition}[Arco parametrizado]\label{def:b2:curves:arc}
Um \emph{arco parametrizado}\index{arco parametrizado} de classe
$\mathcal{C}^k$ ($k \geq 1$) é
uma aplicação $\gamma \colon I \to \R^n$ de classe $\mathcal{C}^k$ num
intervalo $I$. Um ponto $\gamma(t)$ é \emph{regular} se
$\gamma'(t) \neq 0$, e o arco é \emph{regular} se todos os seus
pontos o são. A reta por $\gamma(t)$ dirigida por $\gamma'(t)$
é a \emph{reta tangente}\index{reta tangente} num ponto regular.
\end{definition}

\begin{definition}[Mudança de parâmetro]\label{def:b2:curves:reparam}
Uma \emph{mudança de parâmetro}\index{mudança de parâmetro} de classe
$\mathcal{C}^k$ é um
difeomorfismo $\theta \colon J \to I$ de classe $\mathcal{C}^k$ entre
intervalos ($\theta' \neq 0$ em toda parte). Os arcos $\gamma$ e
$\gamma \circ \theta$ dizem-se \emph{equivalentes}; um
\emph{arco geométrico}\index{arco geométrico} (ou curva) é uma
classe de equivalência. As noções
invariantes por mudança de parâmetro --- a trajetória, a \omterm{def:b2:curves:arc}{reta
tangente}, o \omterm{def:b2:curves:length}{comprimento de arco}, a \omterm{thm:b2:curves:frenet2d}{curvatura} --- são ditas \emph{geométricas}.
\end{definition}

\begin{remark}
A trajetória sozinha não determina o \omterm{def:b2:curves:reparam}{arco geométrico}: as
parametrizações $t \mapsto (\cos t, \sin t)$ em $[0, 2\pi]$ e em
$[0, 4\pi]$ têm a mesma imagem, mas percorrem o círculo uma e
duas vezes. Um \omterm{def:b2:curves:reparam}{arco geométrico} lembra a multiplicidade e a orientação
do percurso, não sua velocidade.
\end{remark}

\begin{example}[A velocidade nada muda de geométrico]\label{ex:b2:curves:speedinvariance}
Parametrize o círculo unitário por
\[
\gamma(t) = (\cos t^2,\ \sin t^2),
\qquad t \in \intcc0{\sqrt{2\pi}} .
\]
A velocidade escalar $\norm{\gamma'(t)} = 2t$ cresce linearmente, e no entanto
\[
L = \int_0^{\sqrt{2\pi}}2t\,\dd t = 2\pi ,
\]
o mesmo \omterm{def:b2:curves:length}{comprimento} que a velocidade constante --- como o
\cref{thm:b2:curves:lengthinv} promete, via a \omterm{def:b2:curves:reparam}{mudança de
parâmetro} $t \mapsto t^2$. A \omterm{def:b2:curves:arc}{reta tangente}, a \omterm{thm:b2:curves:frenet2d}{curvatura}
calculada a partir de \cref{prop:b2:curves:kappaformula} e toda
outra quantidade geométrica também coincidem; só $t = 0$
merece um olhar, onde $\gamma'(0) = 0$ torna \emph{esta
parametrização} irregular embora a trajetória seja um
círculo perfeito. Os enunciados geométricos toleram mal as más
parametrizações: reparametrize primeiro, conclua
depois.
\end{example}

\begin{example}\label{ex:b2:curves:cusp}
O arco $\gamma(t) = (t^2, t^3)$ é $\mathcal{C}^\infty$ mas não é
regular: $\gamma'(0) = (0, 0)$. Sua trajetória, a parábola
semicúbica $y^2 = x^3$, tem uma \emph{cúspide} na origem: a suavidade
da parametrização não impede uma singularidade geométrica
onde a velocidade se anula. É por isso que a hipótese de regularidade
$\gamma' \neq 0$ não é cosmética.
\end{example}

\section{Comprimento de arco}

\begin{definition}[Comprimento de arco]\label{def:b2:curves:length}
Seja $\gamma \colon [a, b] \to \R^n$ um arco $\mathcal{C}^1$. Seu
\emph{comprimento}\index{comprimento de arco} é
\[
L(\gamma) = \int_a^b \norm{\gamma'(t)}\, \dd t ,
\]
em que $\norm{\cdot}$ é a \omterm{def:b2:nvs:norm}{norma} euclidiana. A \emph{função comprimento
de arco}\index{função comprimento de arco} com base em $t_0$ é $s(t) = \int_{t_0}^t
\norm{\gamma'(u)}\,\dd u$.
\end{definition}

\begin{theorem}[O comprimento é geométrico; caracterização
poligonal]\label{thm:b2:curves:lengthinv}
\leavevmode
\begin{enumerate}
\item Se $\theta \colon [c, d] \to [a, b]$ é uma \omterm{def:b2:curves:reparam}{mudança de parâmetro}
$\mathcal{C}^1$, então $L(\gamma \circ \theta) = L(\gamma)$.
\item $L(\gamma)$ é o supremo dos \omterm{def:b2:curves:length}{comprimentos} das
poligonais inscritas:
\[
L(\gamma)
= \sup\Bigl\{\, \sum_{i=1}^{m} \norm{\gamma(t_i) -
\gamma(t_{i-1})} \;:\; a = t_0 < t_1 < \dots < t_m = b \,\Bigr\}.
\]
\end{enumerate}
\end{theorem}

\begin{proof}
\emph{1.} Pela mudança de variáveis $t = \theta(u)$ (volume do primeiro ano de
graduação, válida pois $\theta$ é $\mathcal{C}^1$ monótona),
\[
\int_c^d \norm{(\gamma\circ\theta)'(u)}\,\dd u
= \int_c^d \norm{\gamma'(\theta(u))}\,\abs{\theta'(u)}\,\dd u
= \int_a^b \norm{\gamma'(t)}\,\dd t ,
\]
em que usamos $(\gamma\circ\theta)' = \theta'\cdot
(\gamma'\circ\theta)$ e, se $\theta$ é decrescente, o sinal de
$\theta'$ é absorvido pela inversão dos limites.

\emph{2.} Para qualquer subdivisão, $\gamma(t_i) - \gamma(t_{i-1}) =
\int_{t_{i-1}}^{t_i}\gamma'(t)\,\dd t$, de modo que, pela desigualdade
triangular para integrais, $\norm{\gamma(t_i) - \gamma(t_{i-1})}
\leq \int_{t_{i-1}}^{t_i}\norm{\gamma'}$: toda poligonal é mais curta
que $L(\gamma)$, logo $\sup \leq L(\gamma)$.

Para a desigualdade recíproca, seja $\varepsilon > 0$. Como $\gamma'$
é \omterm{def:b2:metric:continuity}{contínua} no \omterm{def:b2:metric:compact}{compacto} $[a,b]$, ela é uniformemente contínua:
existe $\delta > 0$ com $\norm{\gamma'(t) - \gamma'(u)} \leq
\varepsilon$ sempre que $\abs{t - u} \leq \delta$. Tome uma subdivisão
de passo $\leq \delta$. Em cada pedaço, para $t \in [t_{i-1}, t_i]$,
\[
\gamma(t_i) - \gamma(t_{i-1})
= \int_{t_{i-1}}^{t_i}\gamma'(t)\,\dd t
= (t_i - t_{i-1})\,\gamma'(t_{i-1}) + R_i,
\qquad
\norm{R_i} \leq \varepsilon\,(t_i - t_{i-1}),
\]
pois $R_i = \int_{t_{i-1}}^{t_i}(\gamma'(t) -
\gamma'(t_{i-1}))\,\dd t$. Logo
\[
\norm{\gamma(t_i) - \gamma(t_{i-1})}
\geq (t_i - t_{i-1})\norm{\gamma'(t_{i-1})}
- \varepsilon (t_i - t_{i-1}) .
\]
Somando e comparando $\sum (t_i -
t_{i-1})\norm{\gamma'(t_{i-1})}$ com $\int_a^b\norm{\gamma'}$ (uma
soma de Riemann da função \omterm{def:b2:metric:continuity}{contínua} $\norm{\gamma'}$, a menos de
$\varepsilon(b - a)$ da integral para $\delta$ pequeno o bastante, de novo pela
\omterm{def:b2:metric:continuity}{continuidade} uniforme), obtemos uma poligonal de \omterm{def:b2:curves:length}{comprimento} $\geq
L(\gamma) - 2\varepsilon(b - a)$. Fazendo $\varepsilon \to 0$
prova-se a afirmação.
\end{proof}

\begin{example}[Arquimedes e as poligonais inscritas]\label{ex:b2:curves:archimedes}
Para o círculo unitário, o $n$-ágono regular inscrito tem \omterm{def:b2:curves:length}{comprimento}
$L_n = 2n\sin\frac\pi n$, e o desenvolvimento $\sin x = x -
\frac{x^3}6 + O(x^5)$ dá
\[
L_n = 2\pi - \frac{\pi^3}{3n^2} + O\Bigl(\frac1{n^4}\Bigr) :
\]
os \omterm{def:b2:curves:length}{comprimentos} poligonais do~\cref{thm:b2:curves:lengthinv}
convergem \emph{quadraticamente}. Numericamente: $L_6 = 6$
(o hexágono, dando o rústico $\pi > 3$), enquanto $L_{96} =
192\sin\frac{\pi}{96} \approx 6.28206$ contra $2\pi \approx
6.28319$ --- o erro $0.00113$ concorda com o previsto
$\pi^3/(3\cdot96^2) \approx 0.00112$. É por isso que Arquimedes,
dobrando o hexágono cinco vezes até $96$ lados, conseguiu enquadrar
$\pi$ com três algarismos à mão: cada duplicação divide o
erro por quatro. O supremo da caracterização
poligonal não é apenas atingido no limite; ele é
atingido \emph{rápido}, porque uma curva suave se separa de
suas cordas apenas em segunda ordem.
\end{example}

\begin{theorem}[Parametrização por comprimento de arco]\label{thm:b2:curves:arclength}
Seja $\gamma \colon I \to \R^n$ um arco $\mathcal{C}^k$ regular
($k \geq 1$). A \omterm{def:b2:curves:length}{função comprimento de arco} $s$ é um difeomorfismo
$\mathcal{C}^k$ de $I$ sobre um intervalo $J$, e $\tilde\gamma =
\gamma \circ s^{-1}$ satisfaz $\norm{\tilde\gamma'} = 1$
em toda parte. A menos de translação do parâmetro e de orientação,
essa parametrização por \emph{\omterm{def:b2:curves:length}{comprimento de arco}} (ou de \emph{velocidade unitária}) é
única.
\end{theorem}

\begin{proof}
$s'(t) = \norm{\gamma'(t)} > 0$ e $s'$ é $\mathcal{C}^{k-1}$
(composição da aplicação $\gamma'$ de classe $\mathcal{C}^{k-1}$ com a \omterm{def:b2:nvs:norm}{norma},
suave fora de $0$), de modo que $s$ é $\mathcal{C}^k$, estritamente
crescente, uma bijeção sobre $J = s(I)$, e sua inversa é
$\mathcal{C}^k$ pelo teorema da função inversa em dimensão $1$
(volume do primeiro ano de graduação). Então
\[
\tilde\gamma'(\sigma)
= \frac{\gamma'(t)}{s'(t)}
= \frac{\gamma'(t)}{\norm{\gamma'(t)}},
\qquad t = s^{-1}(\sigma),
\]
um vetor unitário. Se $\hat\gamma = \gamma\circ\theta$ é outra
parametrização de velocidade unitária, então $\abs{\theta'} = 1$, logo $\theta'
= \pm 1$ é constante (\omterm{def:b2:metric:continuity}{continuidade}), isto é, $\theta(u) = \pm u + c$.
\end{proof}

\begin{remark}
O \omterm{def:b2:curves:length}{comprimento de arco} é o parâmetro que separa a geometria da
dinâmica. Uma trajetória pode ser percorrida com qualquer perfil
de velocidade --- a física do movimento --- mas toda questão
invariante por parametrização (forma, dobramento,
osculação) tem um relógio canônico, a distância percorrida.
É por isso que todas as fórmulas de \omterm{thm:b2:curves:frenet2d}{curvatura} abaixo são \emph{definidas}
com velocidade unitária e depois \emph{traduzidas} para parametrizações
arbitrárias pela regra da cadeia: os fatores de tradução
são potências de $v = s'$, e acompanhá-los corretamente é todo
o conteúdo da~\cref{prop:b2:curves:kappaformula}.
\end{remark}

\begin{example}[Círculo e hélice]\label{ex:b2:curves:helix}
Para o círculo $\gamma(t) = (R\cos t, R\sin t)$, $\norm{\gamma'} =
R$, de modo que $s = Rt$ e o \omterm{def:b2:curves:length}{comprimento} de uma volta completa é $2\pi R$. Para a
hélice $\gamma(t) = (a\cos t,\ a\sin t,\ bt)$ com $a > 0$,
$\norm{\gamma'(t)} = \sqrt{a^2 + b^2}$ é constante: a hélice é
percorrida com velocidade constante, e $s = t\sqrt{a^2 + b^2}$.
\end{example}

\begin{example}[Comprimento de arco em coordenadas polares]\label{ex:b2:curves:polarlength}
Uma curva polar $r = r(\theta)$ é o arco $\gamma(\theta) =
(r\cos\theta,\ r\sin\theta)$, com
\[
\gamma'(\theta) = (r'\cos\theta - r\sin\theta,\
r'\sin\theta + r\cos\theta),
\qquad
\norm{\gamma'(\theta)}^2 = r'^2 + r^2
\]
(os termos cruzados se cancelam): o elemento polar de \omterm{def:b2:curves:length}{comprimento} é
$\sqrt{r^2 + r'^2}\,\dd\theta$. Para a cardioide $r = 1 +
\cos\theta$:
\[
r^2 + r'^2 = (1 + \cos\theta)^2 + \sin^2\theta = 2 +
2\cos\theta = 4\cos^2\tfrac\theta2,
\]
e em $\intcc{-\pi}{\pi}$ o cosseno do arco metade é
não negativo, logo
\[
L = \int_{-\pi}^{\pi}2\cos\tfrac\theta2\,\dd\theta
= \Bigl[4\sin\tfrac\theta2\Bigr]_{-\pi}^{\pi}
= 4 - (-4) = 8 :
\]
como o arco de cicloide do~\cref{exo:b2:curves:1}, uma curva
construída a partir de círculos tem \omterm{def:b2:curves:length}{comprimento} racional, sem $\pi$
em lugar algum. A fatoração pelo arco metade é o truque padrão
para \omterm{def:b2:curves:length}{comprimentos} de curvas \omterm{def:b2:structures:generated}{geradas} por círculos; quando ele falha (a
elipse), o \omterm{def:b2:curves:length}{comprimento} é uma função genuinamente nova --- uma
integral elíptica, além das formas fechadas elementares.
\end{example}

\section{Curvatura no plano}

Em toda esta seção, os arcos são $\mathcal{C}^2$ e regulares no
plano euclidiano orientado. Parametrizamos por \omterm{def:b2:curves:length}{comprimento de arco} e
escrevemos $T(s) = \tilde\gamma'(s)$ para a tangente unitária e $N(s)$
para o vetor unitário diretamente ortogonal a $T(s)$ (rotação de
$T$ por $+\pi/2$).

\begin{theorem}[Fórmulas de Frenet no plano]\label{thm:b2:curves:frenet2d}
Seja $\tilde\gamma$ um arco $\mathcal{C}^2$ de velocidade unitária no
plano orientado. Existe uma função \omterm{def:b2:metric:continuity}{contínua} $\kappa$, a
\emph{curvatura}\index{curvatura} (algébrica), tal que
\[
T'(s) = \kappa(s)\, N(s),
\qquad
N'(s) = -\kappa(s)\, T(s).
\]
\end{theorem}

\begin{proof}
Como $\norm{T(s)}^2 = 1$ para todo $s$, derivar o produto
escalar dá $2\langle T'(s), T(s)\rangle = 0$: $T'(s)$ é ortogonal a
$T(s)$, logo colinear com $N(s)$ (dimensão $2$); escreva $T'(s)
= \kappa(s) N(s)$ com $\kappa(s) = \langle T'(s), N(s)\rangle$,
\omterm{def:b2:metric:continuity}{contínua}. Do mesmo modo, $N' \perp N$, logo $N' = \lambda T$; e
derivar $\langle T, N\rangle = 0$ dá $\langle T', N\rangle +
\langle T, N'\rangle = \kappa + \lambda = 0$.
\end{proof}

\begin{definition}\label{def:b2:curves:curvature}
Quando $\kappa(s) \neq 0$, o \emph{raio de
curvatura}\index{raio de curvatura} é $R(s) =
1/\abs{\kappa(s)}$ e o \emph{centro de curvatura}\index{centro de
curvatura} é
$\tilde\gamma(s) + \frac{1}{\kappa(s)} N(s)$; o círculo com esse
centro e raio $R(s)$ é o \emph{círculo osculador}\index{círculo osculador}, a
melhor aproximação circular da curva em $\tilde\gamma(s)$.
\end{definition}

\begin{example}[O círculo osculador da exponencial]\label{ex:b2:curves:oscexp}
Para $y = \eu^x$ no ponto $(0, 1)$: $f'(0) = f''(0) = 1$,
de modo que, pela fórmula do gráfico adiante,
\[
\kappa(0) = \frac{1}{(1 + 1)^{3/2}} = \frac1{2\sqrt2},
\qquad R = 2\sqrt2 .
\]
A tangente unitária é $T = \frac{(1, 1)}{\sqrt2}$, a normal
direta $N = \frac{(-1, 1)}{\sqrt2}$, e o \omterm{def:b2:curves:curvature}{centro de
curvatura} é
\[
(0, 1) + 2\sqrt2\cdot\frac{(-1, 1)}{\sqrt2} = (-2,\ 3) :
\]
o \omterm{def:b2:curves:curvature}{círculo osculador} tem equação $(x + 2)^2 + (y - 3)^2 =
8$. Como verificação da afirmação ``melhor aproximação circular'':
resolver a equação do círculo em $y$ perto de $(0,1)$ e
desenvolver dá $y = 1 + x + \frac{x^2}2 + O(x^3)$ ---
exatamente o desenvolvimento de Taylor de segunda ordem de $\eu^x$. O
\omterm{def:b2:curves:curvature}{círculo osculador} casa valor, inclinação \emph{e} segunda
derivada; um círculo tangente comum casaria apenas os
dois primeiros.
\end{example}

\begin{example}[A evoluta de um círculo é seu centro]\label{ex:b2:curves:circleevolute}
Para o círculo de raio $R$ percorrido no sentido anti-horário,
$\kappa = 1/R$ e $N$ aponta para o centro, de modo que o
\omterm{def:b2:curves:curvature}{centro de curvatura} $\tilde\gamma + \frac1\kappa N$ é o
centro do círculo, para todo $s$: o \omterm{def:b2:curves:curvature}{círculo osculador}
de um círculo é o próprio círculo, e o lugar dos \omterm{def:b2:curves:curvature}{centros
de curvatura} colapsa num ponto. Esse caso degenerado
calibra o~\cref{exo:b2:curves:6}: ali a velocidade da evoluta
é $-\frac{\kappa'}{\kappa^2}N$, que se anula
identicamente precisamente quando $\kappa$ é constante.
\end{example}

\begin{proposition}[Curvatura numa parametrização arbitrária]\label{prop:b2:curves:kappaformula}
Para um arco plano $\gamma(t) = (x(t), y(t))$ regular de classe $\mathcal{C}^2$,
\[
\kappa(t)
= \frac{x'(t)\,y''(t) - y'(t)\,x''(t)}
       {\bigl(x'(t)^2 + y'(t)^2\bigr)^{3/2}} ,
\]
em particular $\kappa = \dfrac{y''}{(1 + y'^2)^{3/2}}$ para um gráfico
$y = f(x)$.
\end{proposition}

\begin{proof}
Escreva $v(t) = \norm{\gamma'(t)} = s'(t)$, de modo que $\gamma' = vT$
(compondo os dados de velocidade unitária com $s$). Derivando,
\[
\gamma'' = v'T + v\,T'\cdot s' = v'T + v^2\kappa N .
\]
Tome agora o \omterm{def:b2:linalg:det}{determinante} (na base canônica orientada) de
$(\gamma', \gamma'')$: como $\det(T, T) = 0$ e $\det(T, N) = 1$,
\[
\det(\gamma', \gamma'') = \det(vT,\ v'T + v^2\kappa N)
= v^3\kappa .
\]
O membro da esquerda é $x'y'' - y'x''$, e $v^3 = (x'^2 +
y'^2)^{3/2}$. O caso do gráfico é a parametrização $t \mapsto (t,
f(t))$.
\end{proof}

\begin{example}[Círculo, reta, parábola]\label{ex:b2:curves:kappaexamples}
Uma reta tem $\kappa = 0$ (e reciprocamente: $T' = 0$ significa $T$
constante, logo $\tilde\gamma(s) = \tilde\gamma(0) + sT$, uma reta). O
círculo de raio $R$ percorrido no sentido anti-horário tem $\kappa = 1/R$:
com $\gamma(t) = (R\cos t, R\sin t)$, a fórmula dá $\kappa =
R^2/R^3$. Para a parábola $y = x^2/2$: $\kappa(x) = 1/(1 +
x^2)^{3/2}$, máxima no vértice --- a parábola é mais fortemente
dobrada onde ela dá a volta.
\end{example}

\begin{remark}[Armadilhas comuns em torno da curvatura]
(i) A \omterm{thm:b2:curves:frenet2d}{curvatura} algébrica de um arco plano muda de sinal quando
se inverte a orientação do arco ou do plano: só
$\abs\kappa$ e $R = 1/\abs\kappa$ são puramente geométricos. Um
círculo percorrido no sentido horário tem $\kappa = -1/R$. (ii) A
fórmula do gráfico $\kappa = f''/(1 + f'^2)^{3/2}$ escolhe
silenciosamente a parametrização por $x$; aplicá-la a uma curva
que não é gráfico perto do ponto (tangente vertical) é o
erro clássico. (iii) Num ponto em que $\gamma' = 0$
nada está definido --- nem $T$ nem $\kappa$ --- e a
trajetória pode realmente quebrar
(\cref{ex:b2:curves:cusp}); verifique sempre a regularidade antes de
derivar a tangente unitária. (iv) No espaço, $\kappa =
\norm{T'} \geq 0$ por convenção: não há sinal a errar,
mas também não há sinal a explorar --- a informação do tipo
inflexão migra para a \omterm{thm:b2:curves:frenet3d}{torção}. (v) Por fim, $\kappa$ é
uma derivada \emph{em relação ao \omterm{def:b2:curves:length}{comprimento de arco}}: para uma
parametrização que não é de velocidade unitária, esquecer o fator $v^3$
em \cref{prop:b2:curves:kappaformula} é o erro mais frequente
na prática.
\end{remark}

\begin{figure}[ht]
\centering
\begin{tikzpicture}[scale=1.6]
  % parabola y = x^2/2
  \draw[->] (-2.2,0) -- (2.4,0) node[below] {$x$};
  \draw[->] (0,-0.4) -- (0,2.4) node[left] {$y$};
  \draw[thick, ocDarkBlue, domain=-2.1:2.1, smooth, samples=60]
    plot (\x, {0.5*\x*\x});
  % osculating circle at vertex: kappa = 1 => R = 1, center (0,1)
  \draw[ocRed, dashed] (0,1) circle (1);
  \fill[ocRed] (0,1) circle (0.03) node[right, font=\small]
    {centro de curvatura};
  \fill (0,0) circle (0.03);
  % tangent and normal at vertex
  \draw[->, thick] (0,0) -- (0.7,0) node[below, font=\small] {$T$};
  \draw[->, thick] (0,0) -- (0,0.7) node[left, font=\small] {$N$};
  % a second point
  \fill (1.2,0.72) circle (0.03);
  \draw[->, thick] (1.2,0.72) -- ++(0.448,0.5376)
    node[right, font=\small] {$T$};
  \draw[->, thick] (1.2,0.72) -- ++(-0.5376,0.448)
    node[above, font=\small] {$N$};
\end{tikzpicture}
\caption{A parábola $y = x^2/2$, seu referencial de Frenet móvel $(T,
N)$ e o \omterm{def:b2:curves:curvature}{círculo osculador} no vértice (raio $1$, pois
$\kappa(0) = 1$). O referencial gira à medida que o ponto se move; a \omterm{thm:b2:curves:frenet2d}{curvatura} é
a taxa desse giro por unidade de \omterm{def:b2:curves:length}{comprimento de arco}.}
\label{fig:b2:curves:osculating}
\end{figure}

\begin{theorem}[A curvatura determina a curva]\label{thm:b2:curves:fundamental}
Seja $\kappa \colon J \to \R$ \omterm{def:b2:metric:continuity}{contínua}. Existe um
arco $\mathcal{C}^2$ de velocidade unitária no plano com \omterm{thm:b2:curves:frenet2d}{curvatura}
$\kappa$, e ele é único a menos de uma isometria direta (rotação
seguida de translação).
\end{theorem}

\begin{proof}
\emph{Existência.} Fixe $s_0 \in J$ e ponha $\varphi(s) =
\int_{s_0}^s \kappa(u)\,\dd u$, depois
\[
\tilde\gamma(s) = \Bigl(\int_{s_0}^s \cos\varphi(u)\,\dd u,\
\int_{s_0}^s \sin\varphi(u)\,\dd u\Bigr).
\]
Então $T(s) = (\cos\varphi(s), \sin\varphi(s))$ é um vetor unitário,
$N(s) = (-\sin\varphi, \cos\varphi)$ e
\[
T'(s) = \varphi'(s)\,(-\sin\varphi, \cos\varphi)
= \kappa(s)\,N(s) :
\]
o arco tem velocidade unitária e \omterm{thm:b2:curves:frenet2d}{curvatura} $\kappa$.

\emph{Unicidade.} Sejam $\gamma_1, \gamma_2$ arcos de velocidade unitária com
a mesma \omterm{thm:b2:curves:frenet2d}{curvatura}. Cada tangente unitária se levanta a uma função
ângulo, por uma construção explícita: veja $T_j$ como o
número complexo $z_j = a_j + \iu b_j$ de módulo $1$, escolha
$\varphi_j(0)$ com $z_j(0) = \eu^{\iu\varphi_j(0)}$ e ponha
\[
\varphi_j(s) = \varphi_j(0) + \int_0^s\det\bigl(T_j,
T_j'\bigr)(u)\,\dd u .
\]
De $\abs{z_j} = 1$: $\operatorname{Re}(\conj{z_j}\,z_j') =
0$, logo $\conj{z_j}\,z_j' = \iu\det(T_j, T_j') =
\iu\,\varphi_j'$, isto é, $z_j' = \iu\varphi_j'z_j$; então
\[
\bigl(z_j\,\eu^{-\iu\varphi_j}\bigr)'
= \eu^{-\iu\varphi_j}\bigl(z_j' - \iu\varphi_j'z_j\bigr) = 0,
\]
de modo que $z_j = \eu^{\iu\varphi_j}$ em toda parte: $T_j =
(\cos\varphi_j, \sin\varphi_j)$ com $\varphi_j$ de classe
$\mathcal C^1$. Além disso, $\det(T_j, T_j') = \det(T_j,
\kappa N_j) = \kappa$, logo $\varphi_j' = \kappa$. Portanto $\varphi_2 = \varphi_1 + c$ para uma
constante $c$: $T_2$ é $T_1$ girado pelo ângulo fixo $c$, de modo que,
integrando, $\gamma_2 = \rho(\gamma_1) + w$, em que $\rho$ é a
rotação de ângulo $c$ e $w$ um vetor constante.
\end{proof}

\begin{remark}
Esse é o protótipo unidimensional de um \emph{teorema fundamental
da geometria}: um conjunto completo de invariantes locais (aqui, uma
função) classifica o objeto a menos de movimento rígido. A
versão tridimensional adiante precisa de dois invariantes.
\end{remark}

\begin{example}[Curvatura constante significa círculo]\label{ex:b2:curves:constantkappa}
Tome $\kappa \equiv \kappa_0 > 0$ na fórmula de existência:
$\varphi(s) = \kappa_0 s$ e
\[
\tilde\gamma(s) = \Bigl(\frac{\sin\kappa_0s}{\kappa_0},\
\frac{1 - \cos\kappa_0s}{\kappa_0}\Bigr) :
\]
o círculo de raio $1/\kappa_0$ centrado em $(0,
1/\kappa_0)$, percorrido com velocidade unitária. Pela metade de
unicidade do teorema, \emph{todo} arco de velocidade unitária de
\omterm{thm:b2:curves:frenet2d}{curvatura} constante $\kappa_0$ é um pedaço de círculo de raio
$1/\kappa_0$ (ou uma reta se $\kappa_0 = 0$) --- a recíproca
do cálculo do~\cref{ex:b2:curves:kappaexamples}, e
o caso plano do~\cref{exo:b2:curves:9}.
\end{example}

\begin{example}[Reconstruir uma curva a partir de sua curvatura]\label{ex:b2:curves:reconstruction}
Que curva de velocidade unitária tem \omterm{def:b2:curves:curvature}{raio de curvatura} $R(s) = 1 +
s^2$? Seguindo a demonstração de existência com $\kappa(s) =
\frac1{1+s^2}$ e $s_0 = 0$: $\varphi(s) = \arctan s$, logo
\[
T(s) = (\cos\arctan s,\ \sin\arctan s)
= \Bigl(\frac{1}{\sqrt{1+s^2}},\
\frac{s}{\sqrt{1+s^2}}\Bigr),
\]
e, integrando,
\[
\tilde\gamma(s) = \Bigl(\ln\bigl(s + \sqrt{1 + s^2}\bigr),\
\sqrt{1 + s^2} - 1\Bigr).
\]
Pondo $x = \ln(s + \sqrt{1+s^2})$, isto é, $s = \sinh x$, a
segunda coordenada é $\cosh x - 1$: a curva é a
\emph{catenária} $y = \cosh x - 1$. Isso fecha o círculo com o
\cref{exo:b2:curves:3}, em que calculamos $R = \cosh^2 x = 1 +
\sinh^2 x = 1 + s^2$ diretamente: o teorema fundamental
garante que a catenária é a \emph{única} curva com esse
perfil de \omterm{thm:b2:curves:frenet2d}{curvatura}, a menos de isometria direta.
\end{example}

\begin{remark}[Onde a curvatura é usada em seguida]
A decomposição $\gamma'' = v'T + v^2\kappa N$ obtida na
demonstração da~\cref{prop:b2:curves:kappaformula} é a
cinemática de todo movimento curvo: aceleração tangencial contra
aceleração centrípeta. A \omterm{thm:b2:curves:frenet2d}{curvatura} volta para superfícies
(\cref{ch:b2:surfaces}) pela \omterm{thm:b2:curves:frenet2d}{curvatura} das curvas nelas
desenhadas, e o cálculo de \omterm{pb:b2:curves:1}{envoltórias} do problema de fim de semana
deste capítulo --- evolutas, cáusticas --- é a ótica geométrica das
frentes de onda. O volume do terceiro ano de graduação retoma o ponto de vista
intrínseco para subvariedades de $\R^n$.
\end{remark}

\section{Referencial de Frenet no espaço}

Seja agora $\tilde\gamma \colon J \to \R^3$ um arco $\mathcal{C}^3$ de velocidade
unitária que seja \emph{birregular}: $T'(s) \neq 0$ para
todo $s$. Então $\kappa(s) = \norm{T'(s)} > 0$ define a
\emph{\omterm{thm:b2:curves:frenet2d}{curvatura}} (sem sinal no espaço: não há orientação normal
preferida), e pomos:
\[
N(s) = \frac{T'(s)}{\kappa(s)}
\quad\text{(normal principal)},
\qquad
B(s) = T(s) \wedge N(s)
\quad\text{(binormal)} ,
\]
de modo que $(T, N, B)$ é um referencial ortonormal direto, o
\emph{referencial de Frenet}. O plano por $\tilde\gamma(s)$ \omterm{def:b2:structures:generated}{gerado} por
$T, N$ é o \emph{plano osculador}.

\begin{theorem}[Fórmulas de Frenet no espaço]\label{thm:b2:curves:frenet3d}
Existe uma função \omterm{def:b2:metric:continuity}{contínua} $\tau$, a \emph{torção}\index{torção}, com
\[
T' = \kappa N, \qquad
N' = -\kappa T + \tau B, \qquad
B' = -\tau N .
\]
\end{theorem}

\begin{proof}
A primeira fórmula é a definição de $N$. Cada um dos vetores
$T, N, B$ tem \omterm{def:b2:nvs:norm}{norma} constante $1$ e eles são dois a dois ortogonais;
derivar as seis relações $\langle X, Y\rangle = \delta_{XY}$
mostra que a matriz de $(T', N', B')$ na base $(T, N, B)$ é
antissimétrica: de fato, $\langle X', Y\rangle + \langle X, Y'\rangle = 0$ e
$\langle X', X\rangle = 0$. Sua entrada $(N, T)$ vale $\langle N', T\rangle =
-\langle N, T'\rangle = -\kappa$, e a entrada da coluna $(T, B)$
vale $\langle T', B\rangle = \kappa\langle N, B\rangle = 0$. Chamar a entrada
livre restante de $\tau = \langle N', B\rangle$ dá exatamente as três fórmulas
exibidas: a antissimetria preenche $\langle B', N\rangle = -\tau$ e
$\langle B', T\rangle = 0$. A \omterm{def:b2:metric:continuity}{continuidade} de $\tau = \langle N', B\rangle$ é clara,
pois $N'$ e $B$ são \omterm{def:b2:metric:continuity}{contínuas} ($\tilde\gamma$ é
$\mathcal{C}^3$, logo $N = T'/\kappa$ é $\mathcal{C}^1$).
\end{proof}

\begin{example}[O vetor de Darboux]\label{ex:b2:curves:darboux}
As três fórmulas de Frenet se comprimem numa só. Ponha $\omega(s) =
\tau\,T + \kappa\,B$ (o \emph{vetor de Darboux}). Usando $B
\wedge T = N$, $T \wedge N = B$, $N \wedge B = T$:
\[
\omega \wedge T = \kappa\,N = T', \qquad
\omega \wedge N = \tau\,B - \kappa\,T = N', \qquad
\omega \wedge B = -\tau\,N = B' :
\]
cada vetor do referencial evolui por $X' = \omega \wedge X$, a
assinatura cinemática de uma \emph{rotação instantânea} de
vetor velocidade angular $\omega$. O referencial gira à taxa
$\norm\omega = \sqrt{\kappa^2 + \tau^2}$ em torno do eixo
móvel $\omega$; a \omterm{thm:b2:curves:frenet2d}{curvatura} é a componente do giro em torno da
binormal, a \omterm{thm:b2:curves:frenet3d}{torção} a componente em torno da tangente. Para
a hélice, $\omega$ é um vetor constante ao longo do eixo do cilindro
--- que é exatamente por que o referencial da hélice precessa
\omterm{def:b2:funcseq:def}{uniformemente}. A antissimetria da matriz de Frenet, explorada no
\cref{exo:b2:curves:7}, é a forma matricial desse único
fato geométrico.
\end{example}

\begin{proposition}[A torção mede a planaridade]\label{prop:b2:curves:torsion}
Um arco birregular está contido num plano se e somente se $\tau
\equiv 0$; nesse caso o plano é o plano osculador (constante).
\end{proposition}

\begin{proof}
Se $\tau \equiv 0$, então $B' = 0$, de modo que $B$ é um vetor unitário
constante $B_0$, e
\[
\frac{\dd}{\dd s}\langle \tilde\gamma(s), B_0\rangle
= \langle T(s), B_0\rangle = 0 :
\]
$\langle \tilde\gamma, B_0\rangle$ é constante, logo o arco está num plano
ortogonal a $B_0$. Reciprocamente, se o arco está num plano $P$,
então $T$ e $T'$ (logo $N$) são paralelos à direção de
$P$ para todo $s$; assim $B = T \wedge N$ é uma das duas normais
\omterm{def:b2:hermitian:adjoint}{unitárias} de $P$ e, sendo \omterm{def:b2:metric:continuity}{contínua}, é constante; então $0 = B'
= -\tau N$ com $N \neq 0$ força $\tau \equiv 0$.
\end{proof}

\begin{example}[Um círculo inclinado tem torção nula]
O arco $\gamma(t) = \bigl(\cos t,\ \tfrac{\sin t}{\sqrt2},\
\tfrac{\sin t}{\sqrt2}\bigr)$ está no plano $y = z$ e
é o círculo unitário desse plano (confira: $\norm{\gamma(t)} =
1$, e a base ortonormal $(1,0,0)$,
$(0, \tfrac1{\sqrt2}, \tfrac1{\sqrt2})$ do plano exibe a parametrização
padrão). Sem cálculo de Frenet algum, a
\cref{prop:b2:curves:torsion} prevê $\tau \equiv 0$, e
a binormal fixa tem de ser a normal unitária do plano $\pm(0,
\tfrac1{\sqrt2}, -\tfrac1{\sqrt2})$. A \omterm{thm:b2:curves:frenet3d}{torção} não mede
estar ``inclinado no espaço''; ela mede \emph{sair} de um
plano. Só o $b$ não nulo da hélice, adiante, produz
\omterm{thm:b2:curves:frenet3d}{torção} genuína.
\end{example}

\begin{example}[A hélice]\label{ex:b2:curves:helixfrenet}
Para a hélice $\gamma(t) = (a\cos t, a\sin t, bt)$, $a > 0$,
calculamos $s = ct$ com $c = \sqrt{a^2 + b^2}$. Então
\[
T = \frac1c(-a\sin t,\ a\cos t,\ b),
\qquad
T' \cdot \frac{\dd t}{\dd s}
= \frac{1}{c^2}(-a\cos t, -a\sin t, 0),
\]
logo $\kappa = a/c^2 = a/(a^2 + b^2)$ e $N = (-\cos t, -\sin t,
0)$: a normal principal aponta horizontalmente para o eixo.
Em seguida, $B = T \wedge N = \frac1c(b\sin t, -b\cos t, a)$, e $B'
\frac{\dd t}{\dd s} = \frac{b}{c^2}(\cos t, \sin t, 0) = -\tau N$
dá
\[
\boxed{\ \kappa = \frac{a}{a^2 + b^2}, \qquad
\tau = \frac{b}{a^2 + b^2}. \ }
\]
Os dois invariantes são constantes --- e pode-se mostrar, reciprocamente, que
as únicas curvas birregulares com $\kappa > 0$ constante e
$\tau$ constante são hélices (círculos quando $\tau = 0$). Note os sinais: $b >
0$ dá uma hélice destra de \omterm{thm:b2:curves:frenet3d}{torção} positiva.
\end{example}

\begin{example}[Curvatura e torção sem comprimento de arco; a
cúbica reversa]\label{ex:b2:curves:twistedcubic}
Reparametrizar por \omterm{def:b2:curves:length}{comprimento de arco} é em geral impossível em forma
fechada, de modo que os invariantes precisam ser extraídos das derivadas
brutas. Escreva $v = \norm{\gamma'} = s'$; então $\gamma' =
vT$ e, como na demonstração da
\cref{prop:b2:curves:kappaformula},
\[
\gamma'' = v'T + v^2\kappa N,
\qquad
\gamma' \wedge \gamma'' = v^3\kappa\,(T \wedge N) =
v^3\kappa\,B .
\]
Tomando \omterm{def:b2:nvs:norm}{normas} ($\kappa \geq 0$ no espaço):
\[
\kappa = \frac{\norm{\gamma' \wedge \gamma''}}{v^3} .
\]
Derivando $\gamma''$ mais uma vez e convertendo $N' =
v(-\kappa T + \tau B)$ (regra da cadeia via $s$), a única
componente em $B$ vem do último termo:
\[
\gamma''' = \bigl(v'' - v^3\kappa^2\bigr)T +
\bigl(v'v\kappa + (v^2\kappa)'\bigr)N + v^3\kappa\tau\,B,
\]
de modo que, emparelhando com $\gamma' \wedge \gamma'' = v^3\kappa B$,
\[
\det(\gamma', \gamma'', \gamma''') = \langle\gamma' \wedge
\gamma'',\ \gamma'''\rangle = v^6\kappa^2\tau,
\qquad\text{i.e.}\qquad
\tau = \frac{\det(\gamma', \gamma'',
\gamma''')}{\norm{\gamma' \wedge \gamma''}^2} .
\]
Aplicação à \emph{cúbica reversa} $\gamma(t) = (t,\
t^2,\ t^3)$ em $t = 0$: $\gamma' = (1, 0, 0)$, $\gamma'' =
(0, 2, 0)$, $\gamma''' = (0, 0, 6)$, logo $v = 1$,
\[
\gamma' \wedge \gamma'' = (0, 0, 2), \qquad \kappa(0) = 2,
\qquad
\det(\gamma', \gamma'', \gamma''') = 12, \qquad \tau(0) =
\frac{12}{4} = 3 .
\]
Lição final: as duas fórmulas são razões em que a velocidade
$v$ se cancela exatamente no grau necessário --- $\kappa$ escala
como uma derivada segunda por unidade de \omterm{def:b2:curves:length}{comprimento}, $\tau$ como o
volume misto de três derivadas por área ao quadrado --- e é
\emph{por isso} que elas são geométricas enquanto o próprio $\gamma''$ não
é.
\end{example}

\begin{remark}[Teorema fundamental para curvas no espaço]
Como no plano, o par $(\kappa, \tau)$ com $\kappa > 0$
determina um arco birregular a menos de isometria direta de $\R^3$: as
fórmulas de Frenet formam um sistema diferencial linear para o referencial
$(T, N, B)$, ao qual se aplica a teoria de Cauchy--Lipschitz do
\cref{ch:b2:diffeq}; a ortonormalidade do referencial solução
é preservada porque a matriz de coeficientes é antissimétrica (mesmo
argumento da matriz de Gram do~\cref{exo:b2:curves:7}), e a curva
é recuperada integrando $T$. Deixamos os detalhes ao
leitor como um exercício substancial mas instrutivo.
\end{remark}

\section{Estudo local: posição em relação à tangente}

\begin{proposition}[Forma local num ponto regular]\label{prop:b2:curves:local}
Seja $\gamma$ um arco plano de classe $\mathcal{C}^k$ em $t_0$, com
$p$ o menor índice tal que $\gamma^{(p)}(t_0) \neq 0$ e $q$ o
menor índice $> p$ com $\gamma^{(q)}(t_0)$ não colinear com
$\gamma^{(p)}(t_0)$ (supondo que ambos existam, $q \leq k$). Na base
$(u, v) = (\gamma^{(p)}(t_0), \gamma^{(q)}(t_0))$ centrada em
$\gamma(t_0)$, Taylor--Young dá as coordenadas
\[
X(t) \sim \frac{(t - t_0)^p}{p!},
\qquad
Y(t) \sim \frac{(t - t_0)^q}{q!} .
\]
O retrato local depende apenas das paridades de $p$ e $q$:
\begin{center}
\begin{tabular}{c c l}
$p$ ímpar, $q$ par & \emph{ponto ordinário} & a curva fica de um só
lado da tangente\\
$p$ ímpar, $q$ ímpar & \emph{ponto de inflexão} & a curva cruza sua
tangente\\
$p$ par, $q$ ímpar & \emph{cúspide de primeira espécie} & os dois ramos
em lados opostos\\
$p$ par, $q$ par & \emph{cúspide de segunda espécie} & os dois
ramos do mesmo lado\\
\end{tabular}
\end{center}
\end{proposition}

\begin{proof}
Taylor--Young na ordem $q$ (a função $\gamma$ é
$\mathcal{C}^q$ perto de $t_0$):
\[
\gamma(t) - \gamma(t_0)
= \sum_{j=p}^{q} \frac{(t-t_0)^j}{j!}\,\gamma^{(j)}(t_0)
+ o\bigl((t-t_0)^q\bigr).
\]
Pela escolha de $p$ e $q$, cada $\gamma^{(j)}(t_0)$ com $p
\leq j < q$ é colinear com $u$; recolhendo as componentes na
base $(u, v)$: $X(t) = \frac{(t-t_0)^p}{p!}(1 + o(1))$ e
$Y(t) = \frac{(t-t_0)^q}{q!}(1 + o(1))$. A tabela de sinais de $X$ e
$Y$ para $t \gtrless t_0$ --- governada exatamente pelas paridades ---
dá os quatro retratos: por exemplo, se $p$ é par, $X > 0$ dos
dois lados (os dois ramos partem na direção $+u$: uma cúspide),
e o lado da \omterm{def:b2:curves:arc}{reta tangente} ($\operatorname{sign} Y$) se inverte
com $q$ ímpar.
\end{proof}

\begin{example}
Para $\gamma(t) = (t^2, t^3)$ em $t_0 = 0$
(\cref{ex:b2:curves:cusp}): $\gamma'' (0)= (2, 0)$, $\gamma'''(0) =
(0, 6)$, logo $p = 2$, $q = 3$: uma cúspide de primeira espécie, o
retrato familiar da parábola semicúbica. Para $\gamma(t) = (t,
t^3)$ em $0$: $p = 1$, $q = 3$: inflexão --- a cúbica cruza
sua tangente.
\end{example}

\begin{remark}[Perspectivas dentro deste volume]
As curvas alimentam os capítulos seguintes de três maneiras. Desenhadas numa
superfície, elas definem seus planos tangentes e sua primeira
forma fundamental (\cref{ch:b2:surfaces}), e seus \omterm{def:b2:curves:length}{comprimentos}
são calculados restringindo a métrica ambiente --- o
capítulo adiante é em grande parte este capítulo relativizado. O
cálculo de \omterm{pb:b2:curves:1}{envoltórias} do problema de fim de semana encontra as integrais
duplas no~\cref{ch:b2:multint}, em que a área do \omterm{pb:b2:curves:1}{astroide}
é recalculada pela fórmula de Green
(\cref{exo:b2:multint:5}) --- uma curva, duas teorias,
respostas coincidentes. E o sistema de Frenet já usou as
equações diferenciais lineares do~\cref{ch:b2:diffeq}
(existência, unicidade e o argumento de preservação da
ortogonalidade do~\cref{exo:b2:curves:7}): o teorema fundamental
das curvas é um teorema de equações diferenciais vestido de
geometria.
\end{remark}

\begin{remark}[Método: executando o estudo local]
Na prática, a classificação é uma rotina de quatro passos.
\emph{Um}, derive em $t_0$ até aparecer a primeira derivada
não nula: seu índice é $p$, seu valor o vetor
$u$. \emph{Dois}, continue derivando até aparecer uma derivada
não colinear com $u$: índice $q$, vetor $v$.
\emph{Três}, leia as paridades $(p, q)$ na tabela.
\emph{Quatro}, desenhe: a curva parte ao longo de $+u$ se $p$ é ímpar
(ao longo de $u$ e depois de volta por $u$ se $p$ é par), do lado
de $v$ ditado pelo sinal de $Y$. Duas advertências. O referencial
$(u, v)$ em geral \emph{não} é ortonormal --- a tabela
descreve posições relativas à \omterm{def:b2:curves:arc}{reta tangente}, não ângulos
nem distâncias, de modo que não leia a \omterm{thm:b2:curves:frenet2d}{curvatura} no desenho.
E derivadas intermediárias colineares com $u$ são permitidas
entre os postos $p$ e $q$ (elas só deslocam o desenvolvimento de
$X$); o que \emph{não} pode acontecer é parar na primeira
derivada não nula e chutar $q = p + 1$: para $\gamma(t)
= (t^2, t^4 + t^5)$ o palpite ingênuo $q = 3$ está errado, pois
$\gamma^{(3)}(0)$ ainda é colinear com $\gamma''(0)$ ---
isso é precisamente o~\cref{exo:b2:curves:5}.
\end{remark}

\section{Exercícios}

\begin{exercise}[$\star$]\label{exo:b2:curves:1}
Calcule o \omterm{def:b2:curves:length}{comprimento} de um arco da cicloide $\gamma(t) = (t -
\sin t,\ 1 - \cos t)$, $t \in [0, 2\pi]$. \emph{(Use $1 - \cos t =
2\sin^2(t/2)$.)}
\end{exercise}

\begin{exercise}[$\star$]\label{exo:b2:curves:2}
Calcule a \omterm{thm:b2:curves:frenet2d}{curvatura} da elipse $\gamma(t) = (a\cos t,\ b\sin
t)$ ($a > b > 0$) e localize os pontos de \omterm{thm:b2:curves:frenet2d}{curvatura}
máxima e mínima.
\end{exercise}

\begin{exercise}[$\star$]\label{exo:b2:curves:3}
Mostre que o \omterm{def:b2:curves:length}{comprimento de arco} do gráfico de $f(x) = \cosh x$ em
$[0, x]$ vale $\sinh x$ e calcule a \omterm{thm:b2:curves:frenet2d}{curvatura} dessa
curva (a \emph{catenária}). Verifique que $R(x) = 1/\kappa(x) =
\cosh^2 x$.
\end{exercise}

\begin{exercise}[$\star\star$]\label{exo:b2:curves:4}
(Espiral logarítmica) Sejam $\gamma(t) = e^{t}(\cos t,\ \sin t)$, $t
\in \R$.
Mostre que o ângulo entre $\gamma(t)$ e $\gamma'(t)$ é
constante, calcule o \omterm{def:b2:curves:length}{comprimento de arco} de $\gamma$ em $(-\infty, 0]$
(finito!) e a \omterm{thm:b2:curves:frenet2d}{curvatura}.
\end{exercise}

\begin{exercise}[$\star\star$]\label{exo:b2:curves:5}
Determine $p$, $q$ e a forma local (ordinário, inflexão,
cúspide) de $\gamma(t) = (t^2,\ t^4 + t^5)$ em $t = 0$, e de
$\gamma(t) = (t^3,\ t^4)$ em $t = 0$.
\end{exercise}

\begin{exercise}[$\star\star$]\label{exo:b2:curves:6}
Seja $\gamma$ um arco plano de velocidade unitária com $\kappa(s) > 0$ para
todo $s$, e seja $c(s) = \gamma(s) + \frac{1}{\kappa(s)}N(s)$
o \omterm{def:b2:curves:curvature}{centro de curvatura} (a curva $c$ é a \emph{evoluta}).
Supondo $\kappa$ de classe $\mathcal{C}^1$, mostre que $c'(s) =
-\frac{\kappa'(s)}{\kappa(s)^2}N(s)$: a evoluta é tangente às
retas normais de $\gamma$.
\end{exercise}

\begin{exercise}[$\star\star\star$]\label{exo:b2:curves:7}
Seja $A(s)$ uma família \omterm{def:b2:metric:continuity}{contínua} de matrizes antissimétricas $3 \times 3$
e $F' = F A$ uma solução matricial com $F(s_0)$
ortogonal. Mostre que $F(s)$ é ortogonal para todo $s$.
\emph{(Derive $G = F F^{\mathsf T}$ e use a unicidade em
Cauchy--Lipschitz.)} Explique a relevância para o sistema de Frenet.
\end{exercise}

\begin{exercise}[$\star\star\star$]\label{exo:b2:curves:8}
(\omterm{thm:b2:curves:frenet2d}{Curvatura} total de uma curva \omterm{def:b2:affine:convex}{convexa} fechada) Seja $\tilde\gamma$ um
arco plano fechado $\mathcal{C}^2$ de velocidade unitária e \omterm{def:b2:curves:length}{comprimento} $L$
(de modo que $\tilde\gamma(s + L) = \tilde\gamma(s)$), percorrido uma vez
no sentido anti-horário. Usando a função ângulo $\varphi$ com $T =
(\cos\varphi, \sin\varphi)$ do
\cref{thm:b2:curves:fundamental}, explique por que $\varphi(L) -
\varphi(0)$ é múltiplo de $2\pi$ e mostre que
$\int_0^L \kappa(s)\,\dd s = \varphi(L) - \varphi(0)$. (Para um
círculo de raio $R$: $\int \kappa = \frac1R \cdot 2\pi R = 2\pi$.
O teorema das tangentes girantes afirma o valor $2\pi$ para
toda curva fechada simples; não se pede que você o demonstre.)
\end{exercise}

\begin{exercise}[$\star\star\star$]\label{exo:b2:curves:9}
Mostre que uma curva espacial birregular com $\kappa > 0$ constante e
$\tau = 0$ é (um arco de) um círculo de raio $1/\kappa$.
\emph{(Use \cref{prop:b2:curves:torsion} e depois mostre que o centro
$\gamma + \frac1\kappa N$ é constante.)}
\end{exercise}

\begin{exercise}[$\star$]\label{exo:b2:curves:10}
Calcule o \omterm{def:b2:curves:length}{comprimento de arco} da parábola $y = x^2/2$ em
$\intcc0a$ e mostre que ele vale
\[
\tfrac12\Bigl(a\sqrt{1 + a^2} + \ln\bigl(a + \sqrt{1 +
a^2}\bigr)\Bigr).
\]
\end{exercise}

\begin{exercise}[$\star\star$]\label{exo:b2:curves:11}
Seja $\gamma$ um arco $\mathcal C^2$ regular em $\R^n$ cujas
\omterm{def:b2:curves:arc}{retas tangentes} passam todas por um ponto fixo $P$. Prove que
a trajetória de $\gamma$ está contida numa reta.
\emph{(Parametrize por \omterm{def:b2:curves:length}{comprimento de arco}, escreva $\gamma(s) +
\lambda(s)T(s) = P$ e derive.)}
\end{exercise}

\begin{exercise}[$\star\star\star$]\label{exo:b2:curves:12}
(Teorema fundamental para curvas no espaço) Sejam $\kappa > 0$ e
$\tau$ funções \omterm{def:b2:metric:continuity}{contínuas} num intervalo $J$. Execute
o programa da observação que segue a
\cref{ex:b2:curves:helixfrenet}: (a) mostre que o sistema
linear $F' = FA(s)$, com $A(s)$ a matriz antissimétrica
de Frenet construída a partir de $\kappa, \tau$ e $F(s_0)$ um referencial
ortonormal direto, tem solução global única, que permanece
um referencial ortonormal direto; (b) construa uma curva birregular
de velocidade unitária com \omterm{thm:b2:curves:frenet2d}{curvatura} $\kappa$ e \omterm{thm:b2:curves:frenet3d}{torção} $\tau$;
(c) demonstre a unicidade a menos de uma isometria direta de $\R^3$.
\end{exercise}

\section{Problema: envoltórias --- o astroide, duas evolutas e
uma cáustica}

\begin{omfigure}
\begin{tikzpicture}[scale=2.2]
  \draw[->] (-1.15,0) -- (1.25,0) node[below] {$x$};
  \draw[->] (0,-1.15) -- (0,1.25) node[left] {$y$};
  \draw[omDef] (1,0) -- (0,0.0) (0.966,0) -- (0,0.259)
    (0.866,0) -- (0,0.5) (0.707,0) -- (0,0.707)
    (0.5,0) -- (0,0.866) (0.259,0) -- (0,0.966);
  \draw[omThm, very thick]
    plot[domain=0:360, samples=120]
    ({cos(\x)*cos(\x)*cos(\x)}, {sin(\x)*sin(\x)*sin(\x)});
\end{tikzpicture}

{\small Uma escada de \omterm{def:b2:curves:length}{comprimento} $1$ escorregando por uma parede (posições
em azul) nunca cruza o \omterm{pb:b2:curves:1}{astroide} $x^{2/3} + y^{2/3} = 1$
(vermelho): o \omterm{pb:b2:curves:1}{astroide} é a \emph{\omterm{pb:b2:curves:1}{envoltória}} da família de
segmentos, tangente a cada um deles.}
\end{omfigure}

\begin{problem}[{Problema de fim de semana --- a máquina das envoltórias e
quatro curvas clássicas}]\label{pb:b2:curves:1}
Uma família de retas a um parâmetro geralmente não cobre o
plano \omterm{def:b2:funcseq:def}{uniformemente}: as retas se acumulam ao longo de uma curva tangente a
todas elas, sua \emph{envoltória}\index{envoltória}. Os raios de luz
tornam as \omterm{pb:b2:curves:1}{envoltórias} visíveis como \emph{cáusticas} --- a curva luminosa
com cúspides numa xícara de café. Este problema constrói a
máquina geral das \omterm{pb:b2:curves:1}{envoltórias} e depois a roda quatro vezes: a escada
que escorrega (\omterm{pb:b2:curves:1}{astroide}), as normais da parábola e da
cicloide (evolutas, com o pêndulo de Huygens no fim) e a
cáustica da xícara de café (\omterm{pb:b2:curves:1}{nefroide}). Em todo o problema, $D_t$ designa a
reta de equação $a(t)\,x + b(t)\,y = c(t)$, em que $a, b, c$
são funções $\mathcal C^2$ com $(a(t), b(t)) \neq (0,0)$,
e $\Delta(t) = a(t)b'(t) - a'(t)b(t)$.

\medskip
\noindent\textbf{Parte I --- A máquina das \omterm{pb:b2:curves:1}{envoltórias}.}
\begin{enumerate}
  \item Suponha $\Delta(t) \neq 0$. Mostre que o
        \emph{sistema característico}
        \[
        \begin{cases} a(t)\,x + b(t)\,y = c(t)\\
        a'(t)\,x + b'(t)\,y = c'(t)\end{cases}
        \]
        tem solução única $E(t) = (x(t), y(t))$, dada por
        $x = \dfrac{cb' - c'b}{\Delta}$, $y = \dfrac{ac' -
        a'c}{\Delta}$.
  \item Suponha além disso que $E$ é $\mathcal C^1$ perto de $t$
        com $E'(t) \neq 0$. Derivando a primeira
        equação do sistema, mostre que $a(t)\,x'(t) +
        b(t)\,y'(t) = 0$ e conclua que a curva $E$
        passa por um ponto de $D_t$ \emph{com a
        direção de} $D_t$: a família é tangente a $E$,
        que se chama sua \emph{\omterm{pb:b2:curves:1}{envoltória}}.
  \item Verificação de bom senso: as \omterm{def:b2:curves:arc}{retas tangentes} da parábola $y =
        x^2/2$ nos pontos $(t, t^2/2)$ são $tx - y =
        t^2/2$. Verifique que a máquina das \omterm{pb:b2:curves:1}{envoltórias} devolve a
        própria parábola.
  \item (\omterm{pb:b2:curves:1}{Envoltória} das normais) Seja $\gamma$ de velocidade unitária
        com $\kappa(s) \neq 0$. A reta normal em
        $\gamma(s)$ é $\{M : \langle M - \gamma(s), T(s)
        \rangle = 0\}$. Mostre que seu sistema característico
        força $\langle M - \gamma(s), N(s)\rangle =
        1/\kappa(s)$, logo que o ponto característico é
        o \omterm{def:b2:curves:curvature}{centro de curvatura}: \emph{a \omterm{pb:b2:curves:1}{envoltória} das
        normais é a evoluta}, recuperando o
        \cref{exo:b2:curves:6}. Confira $\Delta(s) =
        \kappa(s)$.
  \item Duas degenerações. Para o feixe $D_\theta : x\cos
        \theta + y\sin\theta = 0$, mostre que
        o ponto característico é a origem para todo $\theta$
        (a ``\omterm{pb:b2:curves:1}{envoltória}'' colapsa num ponto, e $E' = 0$:
        a questão 2 não se aplica). Para uma família de retas
        paralelas ($a, b$ constante), mostre que $\Delta \equiv 0$ e
        que o sistema característico é em geral
        incompatível: não há \omterm{pb:b2:curves:1}{envoltória}.
\end{enumerate}

\medskip
\noindent\textbf{Parte II --- A escada que escorrega e o
\omterm{pb:b2:curves:1}{astroide}.} Um segmento de \omterm{def:b2:curves:length}{comprimento} $1$ escorrega com uma ponta $P_t =
(\cos t, 0)$ no chão e a outra $Q_t = (0, \sin t)$ na
parede, $t \in \intoo0{\pi/2}$.
\begin{enumerate}[resume]
  \item Mostre que a reta $(P_tQ_t)$ tem equação $x\sin t +
        y\cos t = \sin t\cos t$ e que a máquina das \omterm{pb:b2:curves:1}{envoltórias}
        dá o ponto característico
        \[
        E(t) = (\cos^3 t,\ \sin^3 t) :
        \]
        o \emph{astroide}\index{astroide}, de equação
        implícita $x^{2/3} + y^{2/3} = 1$ (estendida aos
        outros quadrantes por simetria).
  \item Mostre que $E'(0) = 0$ e, usando a classificação
        local (\cref{prop:b2:curves:local}), que
        o \omterm{pb:b2:curves:1}{astroide} tem uma cúspide de primeira espécie em $(1, 0)$
        --- e do mesmo modo em seus quatro pontos sobre os eixos.
  \item Calcule $\norm{E'(t)} = \tfrac32\abs{\sin 2t}$ e
        deduza que o \omterm{def:b2:curves:length}{comprimento} total do \omterm{pb:b2:curves:1}{astroide} é $6$.
  \item Onde a escada toca o \omterm{pb:b2:curves:1}{astroide}? Mostre que $E(t) =
        P_t + \sin^2 t\,(Q_t - P_t)$: o ponto de contato
        divide a escada na razão $\sin^2 t : \cos^2
        t$, varrendo-a de uma ponta à outra à medida que a
        escada escorrega.
  \item Calcule a área encerrada pelo \omterm{pb:b2:curves:1}{astroide}: mostre que
        a área no primeiro quadrante é $3\int_0^{\pi/2}\sin^4
        t\cos^2 t\,\dd t$, avalie a integral por
        linearização ($\sin^2 2t = \tfrac{1 - \cos 4t}2$)
        e conclua que a área total é $3\pi/8$.
\end{enumerate}

\medskip
\noindent\textbf{Parte III --- A evoluta da parábola.}
Seja $\gamma(t) = (t, t^2/2)$.
\begin{enumerate}[resume]
  \item Mostre que a reta normal em $\gamma(t)$ tem equação
        $x + t\,y = t + t^3/2$.
  \item Rode a máquina das \omterm{pb:b2:curves:1}{envoltórias}: mostre que a \omterm{pb:b2:curves:1}{envoltória} das
        normais é
        \[
        E(t) = \Bigl(-t^3,\ 1 + \tfrac32 t^2\Bigr),
        \]
        de equação implícita $x^2 = \tfrac8{27}(y - 1)^3$:
        uma parábola semicúbica.
  \item Confira com a questão 4: calcule o \omterm{def:b2:curves:curvature}{centro de
        curvatura} $\gamma(t) + \frac1{\kappa(t)}N(t)$ a partir de
        $\kappa(t) = (1 + t^2)^{-3/2}$
        (\cref{ex:b2:curves:kappaexamples}) e recupere o
        mesmo ponto.
  \item Mostre que a evoluta tem uma cúspide de primeira espécie em
        $(0, 1)$, o \omterm{def:b2:curves:curvature}{centro de curvatura} no vértice ---
        o ponto em que $\kappa$ é extremal, como prevê
        a fórmula $c' = -\frac{\kappa'}{\kappa^2}N$ do
        \cref{exo:b2:curves:6}.
  \item Quantas normais da parábola passam por um dado
        ponto $(x_0, y_0)$? Mostre que a resposta é governada pela
        cúbica $\tfrac{t^3}2 + (1 - y_0)\,t - x_0 = 0$;
        trate completamente o caso do eixo $x_0 = 0$ (uma normal
        para $y_0 < 1$, três para $y_0 > 1$) e interprete
        a evoluta como a curva de transição.
\end{enumerate}

\medskip
\noindent\textbf{Parte IV --- A cáustica da xícara de café.} Raios
paralelos de direção $(1, 0)$ atingem o interior do espelho
circular $x^2 + y^2 = 1$; o raio que atinge $P_\theta =
(\cos\theta, \sin\theta)$ se reflete segundo a lei da
reflexão.
\begin{enumerate}[resume]
  \item Pela simetria especular em relação à normal (o raio),
        justifique que a direção refletida é $v = u -
        2\langle u, n\rangle n$ com $u = (1,0)$, $n =
        (\cos\theta, \sin\theta)$, e calcule $v =
        -(\cos2\theta, \sin2\theta)$.
  \item Mostre que o raio refletido está na reta
        \[
        x\sin 2\theta - y\cos 2\theta = \sin\theta .
        \]
  \item Rode a máquina das \omterm{pb:b2:curves:1}{envoltórias} ($\Delta = 2$): mostre que
        a cáustica é
        \[
        E(\theta) = \Bigl(\tfrac{3\cos\theta -
        \cos3\theta}4,\ \tfrac{3\sin\theta -
        \sin3\theta}4\Bigr),
        \]
        a \emph{nefroide}\index{nefroide}.
  \item Calcule $E'(\theta) = \tfrac32\sin\theta\,
        (\cos2\theta, \sin2\theta)$; confira que a direção
        tangente é a direção do raio refletido (questão
        16), localize as duas cúspides $(\pm\tfrac12, 0)$ e
        mostre que o raio refletido cruza o eixo $y = 0$
        em $x = \frac1{2\cos\theta}$ --- de modo que raios quase axiais
        focalizam em $x = \tfrac12$: a distância focal $R/2$ de um
        espelho de raio $R$.
  \item Mostre que a \omterm{pb:b2:curves:1}{nefroide} tem \omterm{def:b2:curves:length}{comprimento} total $6$ e que,
        perto de $\theta = 0$,
        \[
        E(\theta) - \bigl(\tfrac12, 0\bigr) =
        \bigl(\tfrac34\theta^2 + o(\theta^2),\ \theta^3 +
        o(\theta^3)\bigr) :
        \]
        uma cúspide de primeira espécie, apontando ao longo do eixo.
  \item Explique num parágrafo por que a cáustica é brilhante:
        por cada ponto logo fora da cáustica passam dois
        raios refletidos, e por cada ponto sobre ela os raios
        estão ``infinitamente concentrados'' (a aplicação $(\theta,
        \text{distância ao longo do raio}) \mapsto \R^2$ tem
        ponto crítico exatamente na \omterm{pb:b2:curves:1}{envoltória}).
\end{enumerate}

\medskip
\noindent\textbf{Parte V --- Huygens: a cicloide é sua própria
evoluta.} Sejam $\gamma(t) = (t - \sin t,\ 1 - \cos t)$, $t \in
\intoo0{2\pi}$, um arco da cicloide.
\begin{enumerate}[resume]
  \item Calcule $\kappa(t) = -\dfrac1{4\sin(t/2)}$ e o
        \omterm{def:b2:curves:curvature}{centro de curvatura}; mostre que a evoluta é
        \[
        c(t) = (t + \sin t,\ \cos t - 1),
        \]
        e que a substituição $t = u + \pi$ a exibe
        como a cicloide original transladada por $(\pi, -2)$:
        \emph{a evoluta de uma cicloide é uma cicloide
        congruente} (Huygens).
  \item Verifique que o \omterm{def:b2:curves:curvature}{raio de curvatura} no ápice $t =
        \pi$ vale $4$, que é metade do \omterm{def:b2:curves:length}{comprimento} $8$ de um
        arco (\cref{exo:b2:curves:1}); localize a cúspide da
        evoluta diretamente abaixo do ápice, à distância $4$.
  \item (A propriedade do fio) Seja $\gamma$ de velocidade unitária com
        $\kappa > 0$, $\kappa$ de classe $\mathcal C^1$ e
        $R = 1/\kappa$ estritamente monótona. Usando $c' = R'N$,
        mostre que o \omterm{def:b2:curves:length}{comprimento de arco} da evoluta entre
        $c(s_0)$ e $c(s_1)$ vale $\abs{R(s_1) - R(s_0)}$.
        Interprete: um fio esticado desenrolado da evoluta,
        de \omterm{def:b2:curves:length}{comprimento} $R(s_0)$ no início, tem sua ponta livre
        traçando a curva original --- de modo que um pêndulo oscilando
        entre duas faces cicloidais do relógio de Huygens
        descreve uma cicloide.
  \item Síntese. A máquina da Parte I produziu o
        \omterm{pb:b2:curves:1}{astroide}, uma parábola semicúbica, uma \omterm{pb:b2:curves:1}{nefroide} e uma
        cicloide. Para cada uma das quatro famílias, diga em uma
        frase onde as hipóteses $\Delta \neq 0$ e
        $E' \neq 0$ valeram ou falharam, e que evento geométrico
        (cúspide, foco, degeneração) cada falha de $E' \neq 0$
        sinalizou. Onde devem aparecer os extremos de \omterm{thm:b2:curves:frenet2d}{curvatura} na
        \omterm{pb:b2:curves:1}{envoltória} das normais, e por quê?
\end{enumerate}
\end{problem}
